\section{Loading images}
\label{sec:loading-images}

\subsection{Dropping or browsing a folder}

Click the drop zone or drag a folder of image files onto it. Every
image in the folder is loaded in natural sort order. The first image
is always treated as the reference frame (frame 1 in the UI,
index 0 internally).

Supported formats: \texttt{.tif}, \texttt{.tiff}, \texttt{.png},
\texttt{.bmp}, \texttt{.jpg}, \texttt{.jpeg}. Grayscale is assumed;
RGB inputs are converted to luminance on load.

\begin{warnbox}
All images in the folder must have the same dimensions. If they
don't, loading stops and an error dialog tells you which image is
the first mismatch.
\end{warnbox}

\subsection{The image list}

After loading, the left sidebar shows three columns:

\begin{center}
\begin{tabular}{lll}
  \toprule
  Column & Width & Content \\
  \midrule
  \# & fixed & 1-based frame number \\
  Filename & flex & file name without extension \\
  Region & 50 px & Need / Edit / Add button \\
  \bottomrule
\end{tabular}
\end{center}

The \emph{Region} column button changes its label based on state:

\begin{description}[leftmargin=1.2cm,style=nextline]
  \item[\textcolor{aldicred}{\textbf{Need}}]
    This frame is a reference frame (per the current workflow type)
    and does not yet have a Region of Interest.
  \item[\textcolor{aldicgreen}{\textbf{Edit}}]
    This frame already has a Region of Interest; click to modify it.
  \item[Add]
    This frame is \emph{not} a required reference frame, but you can
    optionally add one.
\end{description}

Clicking any row's button activates Region of Interest editing for
that specific frame; the canvas switches to show that frame's
reference and the \emph{EDITING REGION OF INTEREST FOR FRAME xx}
banner appears.

\subsection{Navigating frames}

The frame slider at the bottom of the canvas scrubs through the
sequence. Clicking a row in the image list also jumps to that frame.
Keyboard: left/right arrow keys move by one frame when the canvas
has focus.
